\begin{otherlanguage}{ngerman}
	\pdfbookmark[0]{Zusammenfassung}{Zusammenfassung}
	\chapter*{Abstract}
		Der Einsatz von künstlicher Intelligenz (KI) wird in der Informatik und Wirtschaft als nahezu unerschöpfliches Innovationspotential verstanden. Gleichzeitig stößt dieser Ansatz aufgrund seines hohen Energie- und Ressourcenverbrauchs, der schwierigen Nachvollziehbarkeit seiner Funktionsweise und der starken Trend-Dynamik auf reale ethische, gesellschaftliche und praktische Grenzen. Insbesondere der hohe Ressourcen- und Energieverbrauch von KI-Systemen wirft Fragen nach deren Vereinbarkeit mit ökologischen und sozialen Zielen auf. Vor allem dort, wo bestehende Problemstellungen bereits durch energieeffizientere und technisch nachvollziehbare Alternativen gelöst werden.

		Diese Arbeit untersucht im Rahmen meines Praxisprojekts, die Vertretbarkeit von KI-basierten Lösungen in der Datenverarbeitung. Es werden technische Alternativen betrachtet und mögliche ökologische und soziale Auswirkungen analysiert. Ziel der Arbeit ist, Kriterien für verantwortungsbewusste und nachhaltige Gestaltung von datenverarbeitenden Systemen zu formulieren, und die Grenzen des gegenwärtigen KI-Trends abzustecken.
\end{otherlanguage}
